problem:  
Let \(P \to B\) and \(P' \to B\) be dual principal \(S^1\)-bundles over a compact smooth manifold \(B\), equipped with closed flux 3‑forms \(H,H' \in \Omega^3(\cdot)\) satisfying the standard *torsional T‑duality* correspondence:
\[
\pi_!(H)=c_1(P'),\quad \pi'_!(H')=c_1(P),
\]
and such that the pair \((P,H)\leftrightarrow (P',H')\) defines an isomorphism of **exact Courant algebroids**
\[
E_H(P) \;\simeq\; E_{H'}(P')
\]
preserving the de Rham Ševera class.  
Assume further that both \(E_H(P)\) and \(E_{H'}(P')\) admit *T‑duality–related* \(S^1\)-invariant **generalized complex structures** \( \mathcal{J} \) and \( \mathcal{J}' \) in the sense of Cavalcanti–Gualtieri, whose integrability data are compatible under the isomorphism above.

For a generalized complex structure \(\mathcal{J}\) with \(+i\)-eigenbundle \(L\subset E_H(P)\otimes\mathbb{C}\), the infinitesimal deformations are governed by the dg‑Lie (or \(L_\infty\)) algebra of the **holomorphic Lie algebroid complex**
\[
\mathcal{C}_{\mathcal{J}}^\bullet \;=\; \big(\Gamma(\!\wedge^\bullet L^*),d_L,[\,\ ,\,]\big).
\]

> **Problem:**  
> 1. Using the explicit T‑duality isomorphism of Courant algebroids, construct the induced map at the level of Lie algebroid complexes
> \[
> \Phi_T: \mathcal{C}_{\mathcal{J}}^\bullet \;\longrightarrow\; \mathcal{C}_{\mathcal{J}'}^\bullet,
> \]
> and determine whether \(\Phi_T\) is a quasi‑isomorphism when \(H,H'\) represent *integral fluxes with torsion components*.  
> 2. Does the presence of a nontrivial torsion class in \(H^3(P,\mathbb{Z})_{\mathrm{tors}}\) obstruct the existence of such a quasi‑isomorphism, even though T‑duality preserves the de Rham class of the flux?  
> 3. Equivalently, can *torsion T‑duality* fail to preserve the derived deformation theory of invariant generalized complex structures, while remaining an isomorphism of Courant algebroids over the real numbers?

Why is it a "good" problem:  
This formulation anchors the question in rigorously defined and existing mathematical structures—exact Courant algebroids, T‑duality correspondences, and the Lie‑algebroid dg‑Lie complexes controlling generalized complex deformations. The problem targets the **novel torsion phenomena** that may alter derived deformation theories without changing the underlying smooth Courant data.  
It thus lies at the crossroads of **topological T‑duality**, **derived deformation theory**, and **integral cohomological effects**. A solution would clarify whether *discrete torsion* in flux compactifications—normally invisible at the de Rham level—affects the derived moduli space of generalized complex or heterotic geometries.  
This is both conceptually deep and technically meaningful: it asks whether the derived equivalence predicted by T‑duality survives after incorporating integral (torsion) data, thereby probing the integral foundations of dualities in generalized geometry. The question is logically sound, unproven, and elegantly simple in formulation—satisfying all the hallmarks of a “good” cutting‑edge mathematical problem.